The only three-element arithmetic progressions in the value set are $\{1,2,3\}$, $\{2,3,4\}$, $\{3,4,5\}$, and $\{1,3,5\}$. Thus $2$ may not lie between $1,3$ in position; $3$ may not lie between $2,4$ or between $1,5$; and $4$ may not lie between $3,5$.
Separate the values before and after $3$. Each of the pairs $\{1,5\}$ and $\{2,4\}$ must lie entirely on one side of $3$. Hence $3$ is first, third, or last.
If $3$ is first, the other restrictions require $1$ before $2$ and $5$ before $4$. There are $\binom42=6$ interleavings of these two ordered pairs. If $3$ is last, reversal gives another $6$.
If $3$ is third, one of the pairs $\{1,5\}$ and $\{2,4\}$ is before it and the other after it. Either assignment works, with $2!\cdot2!=4$ internal orders. Thus there are $8$ such permutations.
The total is $\boxed{6+6+8=20}$.
